% 创新主张与验证证据应逐一对应。
\section{创新论证与验证计划}
\KeyPoint{\StorySixTitle}{\StorySixQuestion}

\subsection{创新点与对照依据}
\ContentSlot{创新论证}{应从命题需求出发说明拟改进的问题、可比较的基线、团队实际贡献及适用条件。}

\subsection{指标状态与验证方法}
\KeyPoint{\StorySevenTitle}{\StorySevenQuestion}
\noindent\begin{tabularx}{\linewidth}{@{}p{0.18\linewidth}XX@{}}
\toprule
指标类别 & 目标配置 & 口径、条件与方法 \\
\midrule
时延 & \TargetLatency & \LatencyVerificationMethod \\
处理帧率 & \TargetFrameRate & \FrameRateVerificationMethod \\
准确性 & \TargetAccuracy & \AccuracyVerificationMethod \\
功耗 & \TargetPower & \PowerVerificationMethod \\
资源使用 & \TargetResourceUse & \ResourceUseVerificationMethod \\
成本 & \TargetCost & \CostVerificationMethod \\
\bottomrule
\end{tabularx}

\ContentSlot{验证组织与证据}{应在审阅后确定适用指标及定义，说明测试环境、输入数据、基线、测量方法和证据保存方式。指标值、估算值与实测值分开标注；不适用指标应删除或说明原因。}
